\documentclass{ibetest}
\begin{document}
\maketest{End-of-Unit Test --- Elec-S1 Part 1}
         {Foundations of DC circuit theory \,$\cdot$\, sections 1.A to 3.C}
         {45 minutes}{50}

% ---------------------------------------------------------------------------
\exercise[1.A $\cdot$ 1.B $\cdot$ 2.C]{Unit conversions}{3 marks}
Give each result exactly, in scientific notation, without rounding.\\
(a) $u = \qty{0.047}{\kilo\volt}$ in millivolts. \quad
(b) A battery of capacity $Q = \qty{3000}{\milli\ampere\hour}$, in coulombs. \quad
(c) $S = \qty{1.5}{\milli\metre\squared}$ in square metres.
\answer{1,1,1}{3.4cm}{%
  (a) $u = 0.047\times10^{3}\ \mathrm{V} = 47\ \mathrm{V} = 47\times10^{3}\ \mathrm{mV}$
      \qquad $\boxed{u = 4.7\times10^{4}\ \mathrm{mV}}$\\[3pt]
  (b) $Q = 3000\times10^{-3}\ \mathrm{A}\times3600\ \mathrm{s} = 3.000\times3600\
      \mathrm{A\,s}$ \qquad $\boxed{Q = 1.08\times10^{4}\ \mathrm{C}}$\\[3pt]
  (c) $S = 1.5\times(10^{-3}\ \mathrm{m})^2$ \qquad
      $\boxed{S = 1.5\times10^{-6}\ \mathrm{m^2}}$}
\begin{scheme}
\pt{1}{1 pt}{$4.7\times10^{4}\ \mathrm{mV}$.}
\pt{2}{1 pt}{$1.08\times10^{4}\ \mathrm{C}$.}
\pt{3}{1 pt}{$1.5\times10^{-6}\ \mathrm{m^2}$.}
\end{scheme}

% ---------------------------------------------------------------------------
\exercise[1.B]{Electric current --- definition}{5 marks}
Define the electric current $i$ through a cross-section of a conductor. Write down the
relation between $i$, the charge $\dd q$ that crosses the section and the duration $\dd t$.
State the name and the symbol of the SI unit of each of these three quantities, and say which
of these units are SI base units.
\answer{2,2,1}{5.6cm}{%
  The electric current is the \textbf{flow rate of electric charge} through the section: the
  charge that crosses the section per unit time. Its conventional direction is that of the
  positive charges (opposite to the motion of the electrons in a metal).\\[4pt]
  \hspace*{5mm}$\boxed{i = \dfrac{\dd q}{\dd t}}$\\[4pt]
  $i$: \textbf{ampere (A)}; \ $q$: \textbf{coulomb (C)}; \ $t$: \textbf{second (s)}.\\
  The ampere and the second are SI base units. The coulomb is a derived unit:
  $1\ \mathrm{C} = 1\ \mathrm{A\,s}$.}
\begin{scheme}
\pt{1}{2 pts}{Charge crossing the section per unit time (1~pt only for ``a flow of
  electrons'' with no ``per unit time'').}
\pt{2}{2 pts}{$i = \dd q/\dd t$.}
\pt{3}{1 pt}{A, C and s with their exact symbols; ampere (and second) base, coulomb derived.}
\end{scheme}
\begin{tip}
This is Test~1, Exercise~2, in which the unit mark was lost (``C\,s$^{-1}$'' for the current,
``coulombs (s)'' for the charge). The answer expected was ampere (A), coulomb (C), second (s).
\end{tip}

% ---------------------------------------------------------------------------
\exercise[1.A $\cdot$ 1.B]{Measuring instruments}{3 marks}
\question{A voltmeter has its V terminal on $B$ and its COM terminal on $A$. It displays
  $+\qty{3.0}{V}$. Therefore:}
\optlist{0}{$u_{AB} = \qty{3.0}{V}$;}
        {1}{$u_{AB} = \qty{-3.0}{V}$;}
        {0}{$V_A = \qty{3.0}{V}$;}
        {0}{the voltmeter is wrongly connected, so nothing can be concluded.}
\why{the display is $V_{\mathrm{V}} - V_{\mathrm{COM}} = V_B - V_A = u_{BA}$, so
  $u_{AB} = -u_{BA} = \qty{-3.0}{V}$.}
\question{To measure the current through a resistor $R$, an ammeter of internal resistance
  $R_A$ is connected:}
\opttwo{0}{in parallel across $R$, with $R_A \gg R$}{1}{in series with $R$, with $R_A \ll R$}
       {0}{in series with $R$, with $R_A \gg R$}{0}{in parallel across $R$, with $R_A \ll R$}
\why{the current must flow through the ammeter, and $R_A$ must not reduce it.}
\question{A voltmeter of internal resistance $R_V$ is connected across a resistor $R$. The
  measurement is not distorted provided that:}
\opts{0}{$R_V \ll R$}{0}{$R_V = R$}{1}{$R_V \gg R$}{0}{$R_V = 0$}
\why{otherwise the voltmeter diverts part of the current flowing through $R$.}

% ---------------------------------------------------------------------------
\exercise[1.C $\cdot$ 3.A]{Reading a circuit diagram}{4 marks}
\begin{center}
\begin{circuitikz}[line width=0.6pt]
  \draw (0,0) to[battery1, invert, l=$e$] (0,4);
  \draw (0,4) -- (2.4,4) -- (4.8,4) -- (7,4);
  \draw (0,0) -- (2.4,0) -- (4.8,0) -- (7,0);
  \draw (2.4,4) to[R, l_=$R_1$] (2.4,2) to[R, l_=$R_2$] (2.4,0);
  \draw (4.8,4) to[R, l=$R_3$] (4.8,2) to[R, l=$R_4$] (4.8,0);
  \draw (2.4,2) to[R, l=$R_5$] (4.8,2);
  \draw (7,4) to[R, l=$R_6$] (7,0);
  \foreach \p in {(0,0),(2.4,0),(4.8,0),(7,0),(0,4),(2.4,4),(4.8,4),(7,4),(2.4,2),(4.8,2)}
    \node[circ] at \p {};
\end{circuitikz}
\end{center}
Base your answers on the diagram exactly as drawn.
\question{How many nodes?}
\opts{0}{2}{1}{4}{0}{6}{0}{10}
\why{the whole top wire is one node ($e$, $R_1$, $R_3$, $R_6$), the whole bottom wire is
  another ($e$, $R_2$, $R_4$, $R_6$), and the two middle points are nodes ($R_1$, $R_2$,
  $R_5$ and $R_3$, $R_4$, $R_5$). ``6'' counts the drawn junction dots.}
\question{How many branches?}
\opts{0}{5}{0}{6}{1}{7}{0}{9}
\why{each of the 7 components joins two different nodes, so each one is a branch.}
\question{How many independent node equations can be written?}
\opts{1}{3}{0}{4}{0}{6}{0}{7}
\why{$N - 1 = 4 - 1 = 3$: the equation of the last node follows from the others.}
\question{How many independent meshes?}
\opts{0}{3}{1}{4}{0}{5}{0}{7}
\why{$m = B - N + 1 = 7 - 4 + 1 = 4$ (the four windows of the diagram).}

% ---------------------------------------------------------------------------
\exercise[2.A $\cdot$ 2.B]{Characteristics and Ohm's law}{3 marks}
\question{The characteristic $u = g(i)$ of an ohmic conductor is a straight line through the
  origin, of slope $\qty{2.0}{\volt\per\milli\ampere}$. Its conductance is:}
\opts{0}{\qty{2.0}{\milli\siemens}}{1}{\qty{0.50}{\milli\siemens}}{0}{\qty{2.0}{\kilo\siemens}}
     {0}{\qty{500}{\siemens}}
\why{the slope of $u = g(i)$ is $R = \qty{2.0}{V}/\qty{1}{mA} = \qty{2.0}{\kilo\ohm}$, so
  $G = 1/R = \qty{0.50}{\milli\siemens}$.}
\question{In the generator convention, Ohm's law for a resistor reads:}
\opts{0}{$u = Ri$}{1}{$u = -Ri$}{0}{$u = i/R$}{0}{$u = -i/R$}
\why{in the generator convention the arrows of $u$ and $i$ point the same way, so the sign
  changes.}
\question{In the $(u, i)$ plane, a vertical line $u = E$ is the characteristic of:}
\opttwo{1}{an ideal voltage source}{0}{an ideal current source}{0}{an ideal diode}
       {0}{an ohmic conductor}
\why{the voltage is $E$ whatever the current.}

% ---------------------------------------------------------------------------
\exercise[1.B $\cdot$ 2.B $\cdot$ 2.C]{Drift velocity and resistance of a cable}{7 marks}
A copper cable of length $L$ and cross-sectional area $S$ carries a steady current $i$.\\
(a) Determine the drift speed $v$ of the conduction electrons.
(b) How long does an electron take to travel the whole length of the cable? Why does a lamp
at the end of the cable nevertheless light up almost instantly?
(c) Determine the resistance $R$ of the cable, then the power $p$ it dissipates.
\data{$L = \qty{10}{m}$ \hspace{12mm} $S = \qty{2.5}{\milli\metre\squared}$ \hspace{12mm}
      $i = \qty{4.0}{A}$\\[2pt]
      $n = \qty{8.0e28}{\per\metre\cubed}$ \hspace{10mm} $e = \qty{1.6e-19}{C}$
      \hspace{10mm} $\rho = \qty{1.7e-8}{\ohm\metre}$}
\answer{1,1,1,1,1,1,1}{9.2cm}{%
  (a) $i = n\,e\,v\,S$ \ $\Rightarrow$ \ $\boxed{v = \dfrac{i}{n\,e\,S}}$ \qquad
      with \ $S = 2.5\times10^{-6}\ \mathrm{m^2}$\\[3pt]
  \hspace*{5mm}$v = \dfrac{4.0\ \mathrm{A}}{8.0\times10^{28}\ \mathrm{m^{-3}}\times1.6\times
      10^{-19}\ \mathrm{C}\times2.5\times10^{-6}\ \mathrm{m^2}}$\\[3pt]
  \hspace*{5mm}$\phantom{v} = \dfrac{4.0}{8.0\times1.6\times2.5\times10^{28-19-6}}\
      \mathrm{m\,s^{-1}} = \dfrac{4.0}{3.2\times10^{4}}\ \mathrm{m\,s^{-1}}$\\[3pt]
  \hspace*{5mm}$\boxed{v = 1.25\times10^{-4}\ \mathrm{m\,s^{-1}} \approx
      1.3\times10^{-4}\ \mathrm{m\,s^{-1}}}$\\[4pt]
  (b) $\Delta t = \dfrac{L}{v} = \dfrac{10\ \mathrm{m}}{1.25\times10^{-4}\
      \mathrm{m\,s^{-1}}}$ \qquad $\boxed{\Delta t = 8.0\times10^{4}\ \mathrm{s}
      \approx 22\ \mathrm{h}}$\\[2pt]
  \hspace*{5mm}The lamp lights at once because it is the \emph{electric field} (the signal)
  that is set up along the whole circuit almost at the speed of light. All the electrons,
  including those already in the lamp, start drifting together. The drift velocity is not the
  speed of the signal.\\[4pt]
  (c) $\boxed{R = \rho\,\dfrac{L}{S}} = \dfrac{1.7\times10^{-8}\times10}{2.5\times10^{-6}}\
      \Omega = 6.8\times10^{-2}\ \Omega$ \qquad $\boxed{R = 68\ \mathrm{m\Omega}}$\\[3pt]
  \hspace*{5mm}$p = R\,i^2 = 6.8\times10^{-2}\ \Omega\times(4.0\ \mathrm{A})^2
    = 6.8\times10^{-2}\times16\ \mathrm{W}$ \qquad $\boxed{p \approx 1.1\ \mathrm{W}}$}
\begin{scheme}
\pt{1}{1 pt}{Literal expression $v = i/(neS)$, letters only.}
\pt{2}{1 pt}{$S$ in $\mathrm{m^2}$ and numerical expression written with units.}
\pt{3}{1 pt}{$v = 1.3\times10^{-4}\ \mathrm{m\,s^{-1}}$ (the power of ten is the point).}
\pt{4}{1 pt}{$\Delta t = L/v = 8.0\times10^{4}\ \mathrm{s} \approx 22\ \mathrm{h}$.}
\pt{5}{1 pt}{Signal (field) versus drift of the carriers.}
\pt{6}{1 pt}{$R = \rho L/S = 68\ \mathrm{m\Omega}$.}
\pt{7}{1 pt}{$p = Ri^2 \approx 1.1\ \mathrm{W}$.}
\end{scheme}
\begin{tip}
Test~1, Ex.~4 revisited: letters first ($v = i/neS$), then numbers \emph{with units}, then the
powers of ten on their own: $8.0\times1.6\times2.5 = 32$ and $10^{28-19-6} = 10^{3}$.
\end{tip}

% ---------------------------------------------------------------------------
\exercise[2.A $\cdot$ 2.B]{Least-squares estimate of a resistance}{4 marks}
The characteristic of a resistor was measured in the receiver convention:
\data{\begin{tabular}{l|cccc}
      $i_k$ (mA) & 1.0 & 2.0 & 3.0 & 4.0\\ \hline
      $u_k$ (V)  & 2.1 & 4.1 & 5.9 & 8.0
      \end{tabular}}
Compute the least-squares estimate $\hat R = \sum_k u_k i_k \big/ \sum_k i_k^2$, then the
conductance $G$. What power does the resistor receive at the last measurement point?
\answer{1,1,1,1}{5.0cm}{%
  $\sum u_k i_k = 2.1 + 8.2 + 17.7 + 32.0 = 60.0\ \mathrm{V\,mA}$ \qquad
  $\sum i_k^2 = 1 + 4 + 9 + 16 = 30\ \mathrm{mA^2}$\\[3pt]
  $\hat R = \dfrac{60.0}{30}\ \mathrm{k\Omega}$ \qquad
  $\boxed{\hat R = 2.00\ \mathrm{k\Omega}}$ \qquad
  $G = \dfrac{1}{\hat R}$ \qquad
  $\boxed{G = 0.50\ \mathrm{mS} = 5.0\times10^{-4}\ \mathrm{S}}$\\[3pt]
  Last point: $p = u\,i = 8.0\ \mathrm{V}\times4.0\ \mathrm{mA}$ \qquad
  $\boxed{p = 32\ \mathrm{mW}}$ \ (received: receiver convention, $p > 0$)}
\begin{scheme}
\pt{1}{1 pt}{Both sums, with units.}
\pt{2}{1 pt}{$\hat R = 2.00\ \mathrm{k\Omega}$.}
\pt{3}{1 pt}{$G = 0.50\ \mathrm{mS}$.}
\pt{4}{1 pt}{$p = 32\ \mathrm{mW}$ (V$\times$mA $=$ mW).}
\end{scheme}

% ---------------------------------------------------------------------------
\exercise[3.A]{A resistive circuit}{6 marks}
Determine the equivalent resistance $\Req$ seen by the source, the current $i$ it delivers,
the voltage $u$ across $R_2$ (use a voltage divider), then the currents $i_2$ and $i_3$.
Check the node law.
\begin{center}
\begin{circuitikz}[line width=0.6pt]
  \draw (0,0) to[battery1, invert, l=$E$] (0,2.6) to[R, l=$R_1$] (3,2.6) -- (5.2,2.6);
  \draw (3,2.6) to[R, l_=$R_2$] (3,0);
  \draw (5.2,2.6) to[R, l=$R_3$] (5.2,0);
  \draw (0,0) -- (3,0) -- (5.2,0);
  \node[circ] at (3,2.6) {}; \node[circ] at (3,0) {};
  \draw[-{Latex[length=1.8mm]}] (0,2.6) ++(0.25,0) -- ++(0.35,0) node[midway, above] {$i$};
  \draw[-{Latex[length=1.8mm]}] (3,2.45) -- (3,2.15) node[midway, left] {$i_2$};
  \draw[-{Latex[length=1.8mm]}] (5.2,2.45) -- (5.2,2.15) node[midway, left] {$i_3$};
  \draw[-{Latex[length=1.8mm]}] (3.55,0.55) -- (3.55,2.05) node[midway, right] {$u$};
\end{circuitikz}
\end{center}
\data{$E = \qty{12}{V}$ \hspace{10mm} $R_1 = \qty{2.0}{\kilo\ohm}$ \hspace{10mm}
      $R_2 = \qty{6.0}{\kilo\ohm}$ \hspace{10mm} $R_3 = \qty{3.0}{\kilo\ohm}$}
\answer{1,1,1,1,1,1}{7.6cm}{%
  $R_2 \parallel R_3$: \ $R_{23} = \dfrac{R_2R_3}{R_2 + R_3} = \dfrac{18}{9.0}\
   \mathrm{k\Omega} = 2.0\ \mathrm{k\Omega}$ \qquad
  $\Req = R_1 + R_{23}$ \qquad $\boxed{\Req = 4.0\ \mathrm{k\Omega}}$\\[3pt]
  $i = \dfrac{E}{\Req} = \dfrac{12\ \mathrm{V}}{4.0\ \mathrm{k\Omega}}$ \qquad
  $\boxed{i = 3.0\ \mathrm{mA}}$\\[3pt]
  Divider ($R_1$ in series with $R_{23}$):
  $u = \dfrac{R_{23}}{R_1 + R_{23}}\,E = \dfrac{2.0}{4.0}\times12\ \mathrm{V}$ \qquad
  $\boxed{u = 6.0\ \mathrm{V}}$\\[3pt]
  $i_2 = \dfrac{u}{R_2} = \dfrac{6.0\ \mathrm{V}}{6.0\ \mathrm{k\Omega}} = 1.0\ \mathrm{mA}$
  \qquad $i_3 = \dfrac{u}{R_3} = \dfrac{6.0\ \mathrm{V}}{3.0\ \mathrm{k\Omega}}
  = 2.0\ \mathrm{mA}$\\[3pt]
  Node law: $i = i_2 + i_3$: \ $3.0\ \mathrm{mA} = 1.0\ \mathrm{mA} + 2.0\ \mathrm{mA}$.
  \checkmark}
\begin{scheme}
\pt{1}{1 pt}{$R_{23} = 2.0\ \mathrm{k\Omega}$.}
\pt{2}{1 pt}{$\Req = 4.0\ \mathrm{k\Omega}$.}
\pt{3}{1 pt}{$i = 3.0\ \mathrm{mA}$.}
\pt{4}{1 pt}{$u = 6.0\ \mathrm{V}$ with the divider (not with $R_2$ alone).}
\pt{5}{1 pt}{$i_2 = 1.0\ \mathrm{mA}$ and $i_3 = 2.0\ \mathrm{mA}$.}
\pt{6}{1 pt}{Node law checked.}
\end{scheme}
\begin{tip}
The divider formula applies to resistors in \emph{series}: first replace $R_2 \parallel R_3$
by $R_{23}$, then divide between $R_1$ and $R_{23}$.
\end{tip}

% ---------------------------------------------------------------------------
\exercise[3.A]{Millman's theorem with a reversed source}{4 marks}
The ground (bottom wire) is the reference of potentials. Note the orientation of the source
$E_2$. Determine the potential $V_A$ of node $A$.
\begin{center}
\begin{circuitikz}[line width=0.6pt]
  \draw (0,0) to[battery1, invert, l=$E_1$] (0,1.4) to[R, l=$R_1$] (0,2.9);
  \draw (3,0) to[battery1, l=$E_2$] (3,1.4) to[R, l=$R_2$] (3,2.9);
  \draw (6,2.9) to[R, l=$R_3$] (6,0);
  \draw (0,2.9) -- (3,2.9) -- (6,2.9);
  \draw (0,0) -- (3,0) -- (6,0);
  \node[circ] at (3,2.9) {}; \node[circ] at (3,0) {};
  \node[above] at (3,2.9) {$A$};
  \node[font=\scriptsize] at (0.3,0.95) {$+$};
  \node[font=\scriptsize] at (3.3,0.45) {$+$};
  \draw (3,0) node[cground, scale=1.4] {};
\end{circuitikz}
\end{center}
\data{$E_1 = \qty{12}{V}$ \hspace{6mm} $R_1 = \qty{4.0}{\ohm}$ \hspace{6mm}
      $E_2 = \qty{4.0}{V}$ \hspace{6mm} $R_2 = \qty{4.0}{\ohm}$ \hspace{6mm}
      $R_3 = \qty{2.0}{\ohm}$}
\answer{1,1,1,1}{5.4cm}{%
  The $+$ terminal of $E_2$ is on the ground side, so the far end of branch 2 is at the
  potential $-E_2$:\\[3pt]
  \hspace*{5mm}$\boxed{V_A = \dfrac{\dfrac{E_1}{R_1} + \dfrac{-E_2}{R_2} + \dfrac{0}{R_3}}
   {\dfrac{1}{R_1} + \dfrac{1}{R_2} + \dfrac{1}{R_3}}}$
  \qquad Numerator: $\dfrac{12}{4.0} - \dfrac{4.0}{4.0} = 3.0 - 1.0 = 2.0\ \mathrm{A}$\\[4pt]
  Denominator: $\dfrac{1}{4} + \dfrac{1}{4} + \dfrac{1}{2} = 1.0\ \mathrm{S}$ \qquad
  $\boxed{V_A = 2.0\ \mathrm{V}}$ \quad (between $-4$ and $12\ \mathrm{V}$ \checkmark)}
\begin{scheme}
\pt{1}{1 pt}{Millman written with the potential $-E_2$ for the reversed branch.}
\pt{2}{1 pt}{Numerator $= 2.0\ \mathrm{A}$.}
\pt{3}{1 pt}{Denominator $= 1.0\ \mathrm{S}$.}
\pt{4}{1 pt}{$V_A = 2.0\ \mathrm{V}$ (with $+E_2$ one finds $4.0\ \mathrm{V}$: 0 for
  this item).}
\end{scheme}

% ---------------------------------------------------------------------------
\exercise[3.B $\cdot$ 3.C]{Thévenin model and operating point}{11 marks}
\begin{center}
\begin{circuitikz}[line width=0.6pt]
  \draw (0,0) to[battery1, invert, l=$E$] (0,2.4) to[R, l=$R_1$] (2.6,2.4)
        to[R, l=$R_3$] (5.2,2.4) node[ocirc] {};
  \draw (2.6,2.4) to[R, l=$R_2$] (2.6,0);
  \draw (0,0) -- (2.6,0) -- (5.2,0) node[ocirc] {};
  \node[circ] at (2.6,2.4) {}; \node[circ] at (2.6,0) {};
  \node[right] at (5.25,2.4) {$A$}; \node[right] at (5.25,0) {$B$};
\end{circuitikz}
\end{center}
\data{$E = \qty{24}{V}$ \hspace{10mm} $R_1 = \qty{3.0}{\ohm}$ \hspace{10mm}
      $R_2 = \qty{6.0}{\ohm}$ \hspace{10mm} $R_3 = \qty{2.0}{\ohm}$ \hspace{10mm}
      $R_c = \qty{12}{\ohm}$}
(a) Determine the resistance $\RTh$ of the Thévenin model seen from $A$ and $B$.
(b) Determine its electromotive force $\ETh$.
(c) Deduce the current $\INo$ of the equivalent Norton model.
\answer{1,1,1,1,1}{5.6cm}{%
  (a) Sources off ($E$ replaced by a wire): $\RTh = R_3 + \dfrac{R_1R_2}{R_1 + R_2}
      = 2.0 + \dfrac{18}{9.0} = 2.0 + 2.0$ \qquad $\boxed{\RTh = 4.0\ \Omega}$\\[4pt]
  (b) Open circuit: no current in $R_3$, so $u_{AB}$ is the voltage across $R_2$ (divider):
      \\[2pt]
  \hspace*{5mm}$\ETh = \dfrac{R_2}{R_1 + R_2}\,E = \dfrac{6.0}{9.0}\times24\ \mathrm{V}$
  \qquad $\boxed{\ETh = 16\ \mathrm{V}}$\\[4pt]
  (c) $\INo = \dfrac{\ETh}{\RTh} = \dfrac{16\ \mathrm{V}}{4.0\ \Omega}$ \qquad
      $\boxed{\INo = 4.0\ \mathrm{A}}$}
A load $R_c = \qty{12}{\ohm}$ is now connected between $A$ and $B$.\\
(d) Determine the operating point $(u^\star, i^\star)$ analytically. Sketch both
characteristics in the $(u, i)$ plane and locate the operating point.
(e) Compute the efficiency $\eta$ of the transfer.
(f) Which value of $R_c$ would maximise the power received by the load? Give this maximum
power.
\answer[6]{1,1,1,1,1,1}{8.4cm}{%
  \begin{minipage}[t]{0.56\linewidth}\raggedright
  (d) $i^\star = \dfrac{\ETh}{\RTh + R_c} = \dfrac{16}{4.0 + 12}\ \mathrm{A}$ \quad
      $\boxed{i^\star = 1.0\ \mathrm{A}}$\\[3pt]
  \hspace*{5mm}$u^\star = R_c\, i^\star$ \quad $\boxed{u^\star = 12\ \mathrm{V}}$\\[3pt]
  Source: line from $(\ETh, 0) = (16\ \mathrm{V}, 0)$ to
  $(0, \ETh/\RTh) = (0, 4.0\ \mathrm{A})$. Load: line through the origin, slope
  $1/R_c$. They cross at $(12\ \mathrm{V}, 1.0\ \mathrm{A})$.\\[4pt]
  (e) $\eta = \dfrac{R_c}{\RTh + R_c} = \dfrac{12}{16}$ \quad
      $\boxed{\eta = 75\,\%}$
  \end{minipage}\hfill
  \begin{minipage}[t]{0.42\linewidth}\centering\vspace{2pt}
  \begin{tikzpicture}[x=0.19cm, y=0.62cm, >={Latex[length=1.6mm]}, line width=0.5pt]
    \draw[->] (0,0) -- (19,0) node[right] {$u$ (V)};
    \draw[->] (0,0) -- (0,4.8) node[above] {$i$ (A)};
    \foreach \x in {4,8,12,16} \draw (\x,0.08) -- (\x,-0.08) node[below, font=\tiny] {\x};
    \foreach \y in {1,2,3,4} \draw (0.3,\y) -- (-0.3,\y) node[left, font=\tiny] {\y};
    \draw[thick] (16,0) -- (0,4);
    \draw[thick, dashed] (0,0) -- (18,1.5);
    \fill (12,1) circle (2pt);
    \draw[densely dotted] (12,0) -- (12,1);
    \node[font=\scriptsize, anchor=south west] at (3,3.25) {source};
    \node[font=\scriptsize, anchor=north] at (16.5,1.3) {load};
  \end{tikzpicture}
  \end{minipage}\\[5pt]
  (f) Matching: $\boxed{R_c = \RTh = 4.0\ \Omega}$ \qquad
  $P_c^{\max} = \dfrac{\ETh^2}{4\RTh} = \dfrac{256}{16}\ \mathrm{W}$ \qquad
  $\boxed{P_c^{\max} = 16\ \mathrm{W}}$ \quad (with $R_c = 12\ \Omega$: $P_c = 12\
  \mathrm{W}$ only)}
\begin{scheme}
\pt{1}{1 pt}{Sources off, $\RTh = R_3 + R_1 \parallel R_2$.}
\pt{2}{1 pt}{$\RTh = 4.0\ \Omega$.}
\pt{3}{1 pt}{Open circuit, divider across $R_2$ (no current in $R_3$).}
\pt{4}{1 pt}{$\ETh = 16\ \mathrm{V}$.}
\pt{5}{1 pt}{$\INo = 4.0\ \mathrm{A}$.}
\pt{6}{1 pt}{$i^\star = 1.0\ \mathrm{A}$.}
\pt{7}{1 pt}{$u^\star = 12\ \mathrm{V}$.}
\pt{8}{1 pt}{Sketch: both lines with their intercepts, intersection marked.}
\pt{9}{1 pt}{$\eta = 75\,\%$.}
\pt{10}{1 pt}{$R_c = \RTh = 4.0\ \Omega$.}
\pt{11}{1 pt}{$P_c^{\max} = 16\ \mathrm{W}$.}
\end{scheme}

\finish

\ifkey
\par\vspace{4mm}
\begin{tcolorbox}[enhanced, colback=white, colframe=black!60, boxrule=0.5pt, arc=1mm,
  title={\sffamily\bfseries Diagnosis: where did you lose marks?}, fonttitle=\small,
  coltitle=black, colbacktitle=black!10, fontupper=\small]
\begin{tabular}{@{}c l l@{}}
\textbf{Ex.} & \textbf{Lecture section(s)} & \textbf{Redo} \\ \hline
1  & 1.A, 1.B, 2.C: prefixes and units   & Progress tests 1.A, 1.B, 2.C (Ex.~1)\\
2  & 1.B: definition of the current     & Progress test 1.B\\
3  & 1.A, 1.B: voltmeter and ammeter    & Progress tests 1.A (Ex.~3), 1.B (Ex.~3)\\
4  & 1.C, 3.A: nodes, branches, meshes  & Progress tests 1.C (Ex.~3), 3.A (Ex.~3)\\
5  & 2.A, 2.B: characteristics, Ohm     & Progress tests 2.A, 2.B\\
6  & 1.B, 2.B, 2.C: $i = nevS$, Pouillet, Joule & Progress tests 1.B, 2.C (Ex.~4)\\
7  & 2.A: least squares                 & Progress test 2.A (Ex.~4)\\
8  & 3.A: combinations, divider         & Progress test 3.A\\
9  & 3.A: Millman                       & Progress test 3.A (Ex.~4)\\
10 & 3.B, 3.C: Thévenin, operating point & Progress tests 3.B, 3.C
\end{tabular}
\end{tcolorbox}
\fi
\end{document}
